% SEGMENT OLP-0535-S01
% Part: set-theory
% Chapter: story
% Section: urelements

\documentclass[../../../include/open-logic-section]{subfiles}

\begin{document}

\olfileid{sth}{story}{urelements}
\olsection{மூலுறுப்புகள் வேண்டுமா?}

அடுத்த சில அத்தியாயங்களில், திரள்வுறும் படிநிலைமுறைக் கணக்
கருத்தாக்கத்திலிருந்து அடிகோள்களைப் பெற முயல்வோம். ஆனால் மேற்கொண்டு
செல்வதற்கு முன்பு, \emph{மூலுறுப்புகள்} பற்றி இன்னும் சிறிது கூற வேண்டும்.

\olref[approach]{sec} இன் சித்திரம் பழைய \emph{கணங்களிலிருந்து} மட்டுமே
புதிய கணங்களை உருவாக்க அனுமதித்தது. ஆனால் பசுக்கள், பன்றிகள், மணற்துகள்கள்
போன்ற சில \emph{கணமல்லாப் பொருள்களும்} கணங்களின் !!{element}s ஆகலாம்
என்று நாம் அனுமதிக்க விரும்பலாம். அந்நிலையில், கணமல்லாத அடிப்படைப்
பொருள்களாகிய \emph{மூலுறுப்புகளுடன்} தொடங்கி, ஒவ்வொரு படிநிலை~$S$ இலும்
மூலுறுப்புகளையோ $S$ க்கு முந்திய படிநிலைகளில் உருவான கணங்களையோ
(எவ்வகையில் வேண்டுமானாலும் சேர்த்து) கொண்ட “சாத்தியமான எல்லாக்” கணங்களையும்
உருவாக்குவோம் எனக் கூறலாம். விளையும் சித்திரம் இதைப் போன்றிருக்கும்:
\begin{center}
	\begin{tikzpicture}[scale=0.6]
	\tikzset{cut_here/.style={densely dotted}}
	\draw (-4,6) -- (-1, 0) -- (1,0) -- (4,6); 
	\draw[cut_here] (-5,8)--(-4,6);
	\draw[cut_here] (5,8)--(4,6);
	\draw(-1.5, 1)--(1.5, 1);
	\draw(-2, 2)--(2, 2);
	\draw(-2.5, 3)--(2.5,3);
	\draw(-3, 4)--(3, 4);
	\draw(-3.5, 5)--(3.5, 5);
	\draw(-4, 6)--(4, 6);
	\node[label] at (2, 0) {\small 0};
	\node[label] at (2.5, 1) {\small 1};
	\node[label] at (3, 2) {\small 2};
	\node[label] at (3.5, 3) {\small 3};
	\node[label] at (4, 4) {\small 4};
	\node[label] at (4.5, 5) {\small 5};
	\node[label] at (5, 6) {\small 6};
	\end{tikzpicture}
\end{center}
எனவே நாம் ஒரு முடிவெடுக்க வேண்டும்: \emph{மூலுறுப்புகளை அனுமதிக்க வேண்டுமா?}

மெய்யியல் நோக்கில், நமது கோட்பாட்டமைப்பில் மூலுறுப்புகளைச் சேர்ப்பது
பொருத்தமானது. நமது கணக் கோட்பாட்டைப் \emph{பயன்படுத்தத்தக்கதாக} ஆக்குவதே
இதற்கான முதன்மைக் காரணம். இதை விளக்க, இரண்டு கணங்கள் $A$ மற்றும்~$B$
இடையே ஓர் இருபுறச் சார்பு இருந்தால், அப்போதும் அப்போது மட்டுமே, அவை ஒரே
அளவைக் கொண்டவை, அதாவது $\cardeq{A}{B}$, என்று
\olref[sfr][siz][]{chap} இல் கூறியதை நினைவுகூர்க. இப்போது, வயலிலுள்ள
பசுக்களும் தொழுவத்திலுள்ள பன்றிகளும் கணங்களை அமைத்தால், “பசுக்களும்
பன்றிகளும் சம எண்ணிக்கையில் உள்ளன” என்ற கூற்றுக்குக் கணக்கோட்பாட்டு
விளக்கத்தை வழங்கலாம். ஆனால் மூலுறுப்புகளைத் தடைசெய்து, பசுக்களும்
பன்றிகளும் கணங்களை \emph{அமைக்கவில்லை} என்றால், அந்தக் கணக்கோட்பாட்டு விளக்கம்
கிடைக்காது. உண்மையில், கணக் கோட்பாட்டைக் கணங்களையே தவிர வேறு எதற்கும்
பயன்படுத்த நேரடியான வழி நமக்கு இருக்காது. (மூலுறுப்புகளைச் சேர்ப்பதற்கான
மேலும் பல காரணங்களுக்கு, \citealt[pp.~vi, 24, 50--1]{Potter2004} ஐப்
பார்க்கவும்.)

ஆனால் கணிதத்தில் மூலுறுப்புகளை அனுமதிப்பது மிகவும் அரிது. இதற்குக்
கணிசமான ஒரு காரணம், மூலுறுப்புகள் இல்லாமல் கணக் கோட்பாட்டை வகுப்பது
\emph{மிகச் சிறிதளவு} எளிது என்பதாகும். ஆனால் சில நேரங்களில், கணக்
கோட்பாட்டிலிருந்து மூலுறுப்புகளை விலக்குவதற்கு இன்னும் சுவையான
நியாயங்கள் காணப்படுகின்றன:
% SOURCE CORRECTION TA-STH-003: pluralize the generic class of urelements excluded from set theory.
\begin{quote}
	கணக் கோட்பாடே கணிதத்தின் அடித்தளம் என்ற நம்பிக்கைக்கு ஏற்ப, கணங்களைப்
	பற்றி மட்டும் பேசுவதன் மூலம் கணிதம் முழுவதையும் நாம் உள்ளடக்கக்கூடியதாக
	இருக்க வேண்டும்; ஆகவே நமது மாறி, பசுக்கள் மற்றும் பன்றிகள் போன்ற
	பொருள்கள் மீது வீச்சுறக் கூடாது.
	%ஆனால் $C$ ஒரு பசு எனில், $\{C\}$ ஒரு கணம்; ஆனால் அது முறையான கணிதப் பொருள் அல்ல.
	\citep[p.~8]{Kunen1980}
\end{quote}
ஆகவே, பயன்படுதன்மையில் கவனம் செலுத்துவது மூலுறுப்புகளைச்
\emph{சேர்ப்பதைப்} பரிந்துரைக்கும்; குறைப்புவழி அடித்தள நோக்கில்
(கணிதத்தைத் தூய கணக் கோட்பாடாகக் குறைக்கும் நோக்கில்) கவனம் செலுத்துவது
அவற்றை \emph{விலக்குவதைப்} பரிந்துரைக்கலாம். சிறிதளவு சோம்பலும்
மூலுறுப்புகளை விலக்கும் திசையையே சுட்டுகிறது.

நாம் மிகவும் சோம்பலான வழியைப் பின்பற்றுவோம். ஆயினும், இதற்குச் சில
கற்பித்தல்சார் நியாயமும் உண்டு. கணக் கோட்பாட்டின் மெய்யியலைப் படிக்கத்
தொடங்குவதற்கு உங்களுக்குத் தேவையான கணக் கோட்பாட்டுக் கூறுகளை அறிமுகம்
செய்வதே நமது நோக்கம். அந்த மெய்யியல் இலக்கியத்தின் பெரும்பகுதி,
மூலுறுப்புகள் \emph{இல்லாமல்} வகுக்கப்பட்ட கணக் கோட்பாடுகளைப் பேசுகிறது.
எனவே இந்த நூல் அந்த இலக்கியத்தை மிகவும் நெருக்கமாகப் பின்பற்றினால்,
அது அதிகப் பயனுள்ளதாக இருக்கலாம்.

\end{document}
